\documentstyle[% 


righttag% 


,11pt% 


,amssymb% 


,amscd%


,russian%               remove in Eng. 


,izvru%                 Set izven% in Eng. 


]{amsart} 





%% Verify table & complete diagrams; marked \tts





\newcommand{\Sets}{\operatorname{Sets}}


\newcommand{\Id}{\operatorname{Id}}


\newcommand{\In}{\operatorname{in}}


\newcommand{\pr}{\operatorname{pr}}


\newcommand{\Alg}{\operatorname{Alg}}


\newcommand{\Fr}{\operatorname{Fr}}


\newcommand{\R}{\boldsymbol{\frak R}}


\newcommand{\K}{\boldsymbol{\frak K}}


\newcommand{\M}{\boldsymbol{\frak M}}


\newcommand{\T}{\boldsymbol{\frak T}}


\newcommand{\W}{\boldsymbol{\frak W}}


\newcommand{\A}{\frak A}


\newcommand{\B}{\frak B}


\renewcommand{\o}[1]{\overline{#1}}


\renewcommand{\a}[1]{{}^{\ast}{#1}}


\newcommand{\ar}[1]{\overset{#1}{\longrightarrow}{}}


\newcommand{\ep}{\varepsilon}


\newcommand{\epo}{\o{\ep}}


\newcommand{\eo}{\o{e}}


\newcommand{\f}{\varphi}


\newcommand{\p}{\psi}


\newcommand{\te}{\vartheta}


\newcommand{\la}{\lambda}


\newcommand{\tts}[1]{}


 


\theoremstyle{definition} 


\theorembodyfont{\rm} 


\newtheorem{notenonum}{\hskip\parindent Замечание} 


\renewcommand{\thenotenonum}{} 


 


%\hoffset=-15mm         For Eng. 


 


\setcounter{page}{67} 


\god{1998} 


\nomer{1~(428)} %        Remove in Eng. 


%\nomer{Vol.\,42, No.\,1}    Set in Eng. 


% \str{75-81}          Set in Eng. 


\udk{512.57} 


\author{\it С.Н.\,Тронин} 


\title 


{Ретракты и ретракции свободных алгебр} 


 


\date{} 


% \pagestyle{myheadings}     Set in Eng. 


 


\pagestyle{plain} %       Remove in Eng. 


\begin{document} 


 


% \thispagestyle{plain}     Set in Eng. 


 


\maketitle 


 


Данная статья продолжает цикл работ автора \cite{T1}--\cite{T6},


сконцентрированных вокруг вопроса о свободности проективных


конечно порожденных алгебр в многообразиях линейных


алгебр, задаваемых полилинейными тождествами \cite{Dn} (проблемы


2.4 и 2.53).


Проективными алгебрами называются ретракты свободных


алгебр. Когда многообразия не шрейеровы, о свободности


проективных алгебр известно очень мало. Если исключить категории


модулей (которые в принципе тоже можно считать


многообразиями алгебр), и рассматривать ``обычные'' линейные


мультиоператорные алгебры, то проективные алгебры свободны в


многообразиях нильпотентных алгебр \cite{A1}, метабелевых алгебр Ли


\cite{A2}, а также в ряде других, менее "классических",


многообразий \cite{T5}, \cite{T6}. Результаты, аналогичные


доказанным в \cite{A1}, получены также для супералгебр \cite{A3}.





По существу, о произвольной проективной алгебре мало что


известно, кроме того, что существует некоторая определяющая ее


ретракция свободной алгебры. Поэтому естественно считать


исходным объектом исследования сами ретракции свободных алгебр.


Основной результат данной работы состоит в установлении тесной


связи между ретракциями свободных алгебр, определяющими данную


проективную алгебру, и возможными ``претендентами'' на роли изоморфизмов


между данной проективной алгеброй и некоторой свободной алгеброй


(или подходящим родственным объектом). Оказывается, что


конструкция $\f(\pi,\te)$, $\p(\pi,\te)$, изучавшаяся в упомянутых


выше работах автора, дает исчерпывающее описание: любой ``достаточно


хороший'' гомоморфизм (в частности, любой изоморфизм) можно представить


в таком виде при подходящих $\pi,\te$ (теорема 2 и ее следствие).


Теорема 1 описывает ситуацию, когда каждая ретракция определяет


изоморфизм между проективной и свободной алгеброй. Это ---


непосредственное обобщение результата работы \cite{A1}.


В теореме 3 и ее следствиях будут получены необходимые и достаточные условия


того, что получаемые из ретракции гомоморфизмы $\f(\pi,\te)$, $\p(\pi,\te)$


являются изоморфизмами. Некоторые достаточные условия того, что $\f(\pi,\te)$


--- изоморфизм, даны в теоремах 4 и 5.





Приступим к описанию того, как из ретракций могут быть


получены гомоморфизмы. Пусть $\R$ и $\M$


--- некоторые категории,


$\Sets$ --- категория множеств, и пусть дана диаграмма функторов,


коммутативная с точностью до естественных изоморфизмов:


\begin{equation}


\begin{CD}


\R @>{U}>> \M @>{D}>> \R \\


&


   \begin{picture}(0,20)


   \put(25,-5){\vector(-2,1){33}}


   \put(-5,-3){\scriptsize{$\Fr_{\R}$}}


   \put(30,-5){\vector(0,1){20}}


   \put(33,8){\scriptsize{$\Fr_{\M}$}}


   \put(35,-5){\vector(2,1){33}}


   \put(55,-3){\scriptsize{$\Fr_{\R}$}}


   \end{picture} & \\


&& \Sets &&


\end{CD}


\end{equation}


Достаточно даже предполагать только наличие естественного изоморфизма


$\alpha : U\cdot D \cdot \Fr_{\M} @>{\cong}>>\Fr_{\M}$.


По мере необходимости будут вводиться дополнительные условия.


Неформально говоря, $\Fr_{\M}(X)=\Fr(X)$ --- ``свободный'' объект в


$\M$ с ``базисом'' $X$. Под его ретракцией понимается пара морфизмов


$\te : P\to \Fr(X)$, $\pi : \Fr(X)\to P$,


со свойством $\pi\te=1_P$. Рассмотрим следующие морфизмы:


\begin{align*}


\f(\pi,\te):  & P @>{\te}>> \Fr(X) @>{\alpha^{-1}(X)}>> UD\Fr(X) 


@>{UD(\pi)}>> UD(P), \\


\p (\pi,\te) : & UD(P) @>{UD(\te)}>> 


UD\Fr(X) @>{\alpha(X)}>>  \Fr(X) @>{\pi}>> P.


\end{align*}


Это именно та конструкция,о которой шла речь выше.


Морфизмы $\f(\pi,\te)$ и $\p(\pi,\te)$


изучались в упомянутых работах автора, поэтому в данной статье


часть их свойств не будет доказываться. Типичные ситуации,в которых


возникает изображенная выше диаграмма,таковы.





\begin{exam}


Пусть $\M$ --- однородное многообразие линейных


алгебр над коммутативным ассоциативным кольцом с единицей $R$,


$\R$ --- категория $R$-модулей,


$\Fr_{\M}$ и $\Fr_{\R}$ --- функторы взятия свободных алгебр


и модулей соответственно.


Функтор $U$ сопоставляет модулю $M$ аналог тензорной


алгебры модуля в многообразии $\M$. Он заведомо существует, если


$M$ проективен.Фактически вместо всего многообразия и всей категории


модулей достаточно взять только подкатегории проективных объектов.


Функтор $D$ сопоставляет алгебре ее фактор по


``квадрату'' --- идеалу, порожденному всеми словами, в запись которых


входит не менее двух элементов алгебры. В частности, этот функтор


сопоставляет свободной алгебре первую компоненту градуировки.


\end{exam}





\begin{exam}


Все конструкции и результаты сохранятся, если вместо свободных


алгебр однородного многообразия брать их пополнения относительно


стандартной градуировки.


\end{exam}





\begin{exam}


Подобную же диаграмму функторов можно получить в случае,


когда $\M$ --- категория модулей над кольцом $K$ --- полупрямым


произведением $R$ и некоторого идеала ${\frak A}$. Тогда


$U(M)=M\otimes_R K$, $D(N)=N\otimes_K R \cong N/N{\frak A}$.


\end{exam}





Чтобы иметь возможность проводить рассуждения сразу для всех этих


случаев, воспользуемся языком финитарных алгебраических теорий в


смысле Ловера (\cite{Law}, \cite{Shub}, гл.\,18;


\cite{KR}, \cite{BF},


\cite{M}, \cite{W1}, \cite{W2}) (далее называемых


просто теориями). Алгебраическая теория $\T$ --- это категория с


объектами $0,1,2,\dotsc$, причем объект $n+m$ есть декартово


произведение объектов $n$ и $m$, что равносильно существованию


естественного изоморфизма


${\T}(k,n+m) \cong {\T}(k,n) \times {\T}(k,m)$.


Простейший пример --- категория ${\bold{Mat}}(R)$, дуальная к категории


свободных


правых $R$-модулей конечного ранга: множества ${\bold{Mat}}(R)(n,m)$ есть


множества всех $n\times m$-матриц над $R$ (которое уже не обязательно


считать коммутативным). Морфизм алгебраических теорий


$\Upsilon :{\T} \to {\W}$


есть функтор, тождественный на объектах и сохраняющий произведения.


${\T}$-алгебрами называются функторы из ${\T}$ в $\Sets$,


сохраняющие произведения, гомоморфизмами --- естественные


преобразования. Через $\Alg({\T})$ обозначается соответствующая


категория. Морфизму теорий соответствует функтор


$\Upsilon_{\#} : \Alg({\T}) \to \Alg({\W})$,


обладающий правым сопряженным.





Связь с обычными многообразиями универсальных алгебр такова.


Если $\M$ --- многообразие, то по нему однозначно строится


алгебраическая теория: категория $\T$, двойственная к категории


свободных алгебр конечного ранга многообразия $\M$. Произвольной


алгебре $A$ из $\M$ сопоставляется функтор из $\T$ в $\Sets$,


отображающий объект $n$ в $A^n$ --- прямое произведение $n$ экземпляров


$A$. Обратный функтор сопоставляет функтору $F: {\T} \to \Sets$


множество $F(1)$. Если ${\T}$ --- произвольная теория, то


по ней можно построить многообразие следующим образом. В качестве


множества всех $n$-арных операций берется множество ${\T}(n,1)$,


а тождества соответствуют правилам композиции морфизмов ${\T}$.


В конечном счете многообразие определяется алгебраической теорией


с точностью до рациональной (или категорной) эквивалентности.


Если $\Fr_{\M}(X)$ --- свободная алгебра многообразия $\M$


с базисом $X$, а ${\T}$ --- соответствующая теория, то


$\Fr_{\M}(x_1,\dotsc,x_n)={\T}(n,1)$,


а ${\T}(n,m)$ интерпретируется как множество всех гомоморфизмов


из $\Fr_{\M}(x_1,\dotsc,x_m)$ в $\Fr_{\M}(x_1,\dotsc,x_n)$.





Определим $R$-линейную алгебраическую теорию как теорию ${\T}$


вместе с морфизмом теорий


$\Upsilon : {\bold{Mat}}(R) \to {\T}$.


Будем также предполагать, что множества ${\T}(n,0)$ состоят из


одного элемента. Тем самым на множествах ${\T}(n,m)$ вводятся


структуры правых $R$-модулей, причем композиция морфизмов не билинейна.


Выполнены только соотношения:


$0 \cdot f = 0$, $f\cdot (g \pm h)= f\cdot g \pm f\cdot h$.





Далее будет дополнительно предполагаться, что $\Upsilon$ ---


расщепляющийся мономорфизм, т.е. существует морфизм теорий


$\Delta : {\T}\to {\bold{Mat}}(R)$


такой, что $\Delta\cdot\Upsilon = 1_{{\mathrm{Mat}}(R)}$.





В примерах 1, 2, 3 мы имеем дело именно с такими $R$-линейными теориями.


Диаграмма (1) возникает теперь, если взять


$\R=Mod-R=\Alg({\bold{Mat}}(R))$, $\M=\Alg({\T})$,


$D=\Delta_{\#} $, $U=\Upsilon_{\#}$.


Вариант: $\R={\bold P}(R)$, $\M={\bold P}({\T}) $ ---


категории проективных $R$-модулей и проективных ${\T}$-алгебр.





При сделанных предположениях свободные ${\T}$-алгебры


$\Fr(X)=\Fr_{\M}(X)$ обладают следующей структурой:


$\Fr(X)=\Fr^1(X)\oplus\Fr^{\bullet}(X)$,


где $\Fr^1(X)$ --- свободный правый $R$-модуль с базисом $X$,


$\Fr^1(X)\cong D(\Fr(X))$,


$\Fr^{\bullet}(X)$ --- подмодуль, такой что любой гомоморфизм


из $\Fr(X)$ в $\Fr(Y)$ (который к тому же является гомоморфизмом


правых $R$-модулей) отображает $\Fr^{\bullet}(X)$ в


$\Fr^{\bullet}(Y)$.


Заметим, для любого ретракта $P$ объекта $\Fr(X)$


определен подмодуль $P^{\bullet}$ с аналогичным свойством,


выделяющийся прямым слагаемым (прямое


дополнение, определено не однозначно, но $P/P^{\bullet}\cong D(P)$), и


обладающий аналогичным свойством.


Мы будем записывать


элементы $\Fr(X)$ в виде $f=f^1+f^{\bullet}$ или


$f(X)=f^1(X)+f^{\bullet}(X)$, $f^1\in\Fr^1(X)$, $f^{\bullet}\in


\Fr^{\bullet}(X)$,


причем $f^1(X)$ можно мыслить как ``линейную часть'' $f(X)$,


а $f^{\bullet}(X)$ --- как ``сумму членов степеней,больших единицы''.


Поскольку множества элементов $\Fr(X)$ соответствуют гомоморфизмам


свободных алгебр в $\Fr(X)$, то возможны подстановки произвольных


элементов вместо переменных, обладающие обычными свойствами.


Будут использованы векторные обозначения: если даны гомоморфизмы


$


\Fr(X)@>{\eta}>> \Fr(Y)


@>{\mu}>> P $,


$X=\{x_1,\dotsc,x_n\}$, $Y=\{y_1,\dotsc,y_m\},


$


и $\mu(y_j)=p_j$, $1\le j\le m$, $\eta(x_k)=g_k=g_k(Y)$, $1\le k\le n$, то


полагаем $p=\{p_j|1\le j\le m\}$, и $\eta(\mu(x_k))=g_k(p)$, $1\le k\le n$.


Множество $\{g_k(p)|1\le k\le n\}$ будем обозначать через $g(p)$.


В частности, $X=\{x_k|1\le k \le n\}$, $Y=\{y_j|1\le j\le m\}$


укладываются в эту же схему.





Введем следующие обозначения. Пусть дана ретракция


$\pi :\Fr(X) \to P $,


$\te : P \to \Fr(X)$, $\pi\cdot \te=1_P $,


определяющая данную проективную алгебру $P$. Положим


${}^0 P=UD(P)$, ${}^0\pi=UD(\pi)$, ${}^0\te=UD(\te)$,


${}^\ast\pi={}^0\pi\cdot \alpha^{-1}(X):


\Fr(X) \to {}^0 P$,


${}^\ast\te= \alpha(X)\cdot {}^0\te:


{}^0 P \to \Fr(X)$.


Тогда ${}^\ast\pi\cdot {}^\ast\te = 1_{UD(P)}$,


$\f(\pi,\te)= {}^\ast\pi\cdot\te : P


\to {}^0 P $,


$\p(\pi,\te)=\pi\cdot{}^{\ast}\te {}^0 : P \to P$.


Отметим, что обозначения типа ${}^0 P$


введены здесь по аналогии с работой \cite{BCW}.





Пусть $X=\{x_1,x_2,\dotsc,x_n\}$. Предположение о конечности


$X$ не будет играть в дальнейшем почти никакой роли и сделано


лишь для упрощения обозначений. Положим


$f_i=\te\cdot\pi(x_i)$, $i=1,\dotsc,n$,


$f_i=f^1_i+f^{\bullet}_i$, $f=\{f_i\mid 1\le i\le n\}$,


$f^1=\{f^1_i \mid 1\le i\le n\}$.


Тогда из $\pi\cdot\te=1$ и из сделанных выше предположений


будут следовать соотношения:


$ f_i(X)=f_i(f)$, $f^1_i(X)=f^1_i(f^1)$, $1\le i\le n$,


и кроме того, для всех $i\;$


$f^1_i(X)={}^{\ast}\te({}^\ast\pi(x_i))$.





Определим гомоморфизмы


$ \sigma, \tau : \Fr(X) \to \Fr(X)$


для фиксированной ретракции $(\pi,\te)$,


полагая


$$


\tau(x_i) = x_i-f_i+f^1_i = x_i-f^{\bullet}_i \quad, 


\sigma(x_i) = x_i+f_i-f^1_i = x_i+f^{\bullet}_i.


$$





\begin{lem}


Имеет место коммутативная диаграмма


$$


\begin{CD}


\Fr(X) @>{{}^\ast\pi}>> {}^0 P @>{{}^\ast\vartheta}>> \Fr(X) @>{\pi}>> P \\


@VV{\tau}V @VV{\p}V @VV{\sigma}V @VV{\f}V\\


\Fr(X) @>{\pi}>> {}^0 P  @>{\vartheta}>> \Fr(X)  @>{{}^\ast\pi}>> {}^0 P,


\end{CD}


$$


где $\f=\f(\pi,\vartheta)$, $\p=\p(\pi,\vartheta)$.


Если $\f$ --- инъекция, и существует $\sigma^{-1}$, то


существует и $\f^{-1}=\pi\cdot\sigma^{-1}\cdot {}^{\ast}{\te}$. Если


$\p$ --- инъекция, и существует $\tau^{-1}$, то


существует и $\p^{-1}={}^{\ast}\pi\cdot\tau^{-1}\cdot\te $.


\end{lem}





\begin{pf} Положим


$p_i=\pi(x_i)$, ${}^0 p_i={}^{\ast}\pi(x_i)$,


$1\le i\le n$.


Тогда из определений следует,что


$\psi({}^0 p_i)=f^1_i(p)$, $ \varphi(p_i)=f_i({}^0 p)$,


причем $p=f(p)$, ${}^0 p= f^1({}^0 p)$.


Коммутативность достаточно проверить,сравнивая значения гомоморфизмов


на $x_i$ или ${}^0 p_i$. Это достигается легким вычислением,


с учетом того, что подстановки ведут себя обычным образом. Например,


для левого квадрата


$\psi({}^{\ast}\pi(x_i))=f^1_i(p)$,


$\pi(\tau(x_i))=\pi(x_i-f_i+f^1_i)= p_i-f_i(p)+f^1_i(p)$.


Проверим, что указанные выражения действительно определяют обратные для


$\f$ и $\p$. Из $\f\cdot\pi={}^{\ast}\pi\cdot\sigma$


следует


$\f\cdot (\pi\cdot\sigma^{-1}\cdot{}^{\ast}\te)=1$,


следовательно, $\f$ --- сюръекция, а т.к. по условию это


еще и инъекция, то получаем биективность. Аналогично рассуждаем для


случая $\p$.


\end{pf}





Назовем многообразие $\M$ квазиоднородным, если любой эндоморфизм вида


$$


\lambda : \Fr(X)\to \Fr(X),\quad


\lambda(x_i)=x_i+ q_i,\quad q_i \in \Fr^{\bullet}(X)


$$


является мономорфизмом, и квазиполным, если все такие $\lambda$


сюръективны. В примере 1 $\M$ квазиоднородно, в примере 2 ---


квазиоднородно и квазиполно. В примере 3 все зависит от идеала


$\frak A$. Если $\frak A$-адическая фильтрация отделима,


то $\M$ квазиоднородно, а если кольцо $R$ полно относительно


соответствующей топологии, то $\M$ квазиполно.


Во всех этих примерах $\M$ обладает


свойством, которое будет называться квазихопфовостью:


сюръективный эндоморфизм $\Fr(X)$ указанного выше вида


является изоморфизмом (напомним, что хопфовость


\cite{BA} означает, что если $X$ конечно,то любой сюръективный


эндоморфизм $ \Fr(X)$ является изоморфизмом).





\begin{thm}


Если многообразие $\M$ квазиоднородно,то любой морфизм


$\p(\pi,\te)$ инъективен, а если квазиполно, то все


$\p(\pi,\te)$ и $\f(\pi,\te)$ сюръективны. Если $\M$


квазиполно и обладает свойством квазихопфовости, то все $\p(\pi,\te)$


являются биекциями, и в этом случае


$ \Fr(X) \cong \te(P) \coprod Q$,


где $Q$ --- ретракт $\Fr(X)$, являющийся образом идемпотентного


эндоморфизма $\Fr(X)$, отображающего $x_i$ в $x_i-f_i^1$,


$1 \le i \le n$.


\end{thm}





\begin{pf} Утверждения об инъективности


и сюръективности следуют из леммы 1. Если $\M$ квазиполно и


обладает свойством квазихопфовости, то $ \sigma $ ---


биекция. Отсюда следует, что $\psi$ --- также биекция. Заметим, что


$f^1_i(X)=\sum\limits_{1\le j \le n}x_j a_{ji}$, $1\le i \le n$,


где $(a_{ji})$ --- идемпотентная матрица над $R$. Следовательно,


соответствие


$x_i \longmapsto x_i -f^1_i$


определяет идемпотентный эндоморфизм $\Fr(X)$,


образ которого $Q$


есть ретракт $\Fr(X) $. Легко убедиться, что


$Q \cong {}^0 Q= UD(Q)$,


причем $D(Q)$ --- конечно порожденный проективный $R$--модуль


--- образ эндоморфизма $R^n$, задаваемого матрицей


$I_n -(a_{ij})$ (где $I_n$ --- единичная $n\times n$ --матрица).


В свою очередь, $D(P)$


есть образ эндоморфизма с матрицей $(a_{ij}) $.


Теперь рассмотрим цепочку изоморфизмов


$$


\te(P)\coprod Q \cong {}^0 P \coprod {}^0 Q =


UD(P)\coprod UD(Q)


\cong U(D(P)\oplus D(Q)) \cong U(R^n) \cong \Fr(X).


$$


Здесь используются свойства функторов $U$ и $D$,


в частности, их перестановочность с копроизведениями.


Поскольку $\te(P)$ и $Q$


естественным образом вложены в $\Fr(X)$, то существует


гомоморфизм


$\zeta: \te(P) \coprod Q \to \Fr(X)$.


Взяв композицию $\zeta$ с предшествующим изоморфизмом между


$\Fr(X)$ и $\te(P) \coprod Q$, и прослеживая шаг за шагом


действие полученного


гомоморфизма $\Fr(X)\to \Fr(X)$


на элементы $X$, можно легко убедиться, что каждый $x_i$


отображается в элемент вида $x_i+q_i$,


$q_i\in \Fr^{\bullet}(X)$.


Следовательно, применимы условия квазиполноты и квазихопфовости.


\end{pf}





Заметим, что, как показано в \cite{T2}, \cite{T4}, \cite{T6},


в примерах 1 и 2, когда многообразие задается однородными


тождествами (а также в примере 3, если фильтрация отделима),


всегда инъективны не только $\psi $, но и $\varphi$.


В рамках данной статьи нет необходимости в соответствующем


усилении теоремы~1, которое потребовало бы дополнительных условий.


Заметим еще раз, что последнее утверждение


этой теоремы можно рассматривать как обобщение основного


результата работы \cite{A1}.





Теперь настало время вспомнить, что из способа построения диаграммы


(1) следует факт существования естественного изоморфизма


$\beta : \Id_{\R} \to D\cdot U$


такого, что гомоморфизмы\linebreak


$D\Fr(X) @>{\beta(D\Fr(X))}>> DUD\Fr(X)$,


$DUD\Fr(X) @>{D(\alpha(X))}>> D\Fr(X)$


являются взаимно обратными изоморфизмами. Легко проверить, что


$D(\f(\pi,\te))=\beta(D(P))$,


$D(\p(\pi,\te))=\beta(D(P))^{-1}$.


Произвольные гомоморфизмы вида


$\p': {}^0 P\to P$, $\f': P\to {}^0 P$


такие, что


$D(\f')=\beta(D(P))$,


$D(\p')=\beta(D(P))^{-1}$,


будем называть квазиинъективными.


Во всех дальнейших результатах не


предполагается ни квазиоднородности, ни квазиполноты, ни конечности


базисов свободных алгебр.





\begin{lem}


Пусть даны ретракции


$\pi: \Fr(X) \to P$,


$\te: P\to \Fr(X)$,


$\pi\te=1_P$,


$\rho: P \to Q$,


$\lambda: Q\to P$,


$\rho, \lambda=1_Q$.


Тогда


$\f(\rho\cdot\pi,\te\cdot\lambda)=UD(\rho)\cdot\f(\pi,\te)


\cdot\lambda$,


$\p(\rho\cdot\pi,\te\cdot\lambda)=\rho\cdot\p(\pi,\te)


\cdot UD(\lambda )$.


\end{lem}





Доказательство следует из определений.





\begin{lem}


Пусть


$\p': {}^0 P\to P$, $\f': P\to {}^0 P$,


квазиинъективны, и пусть даны любые


$\rho : \Fr(X)\to P$, $\lambda : P\to \Fr(X)$,


$\rho\cdot \lambda=1_P$.


Тогда существуют квазиинъективные


$\p'': {}^0 \Fr(X)\to \Fr(X)$,


$\f'': \Fr(X)\to {}^0 \Fr(X)$


такие, что коммутативны диаграммы


$$


\begin{CD}


P & @>{\f'}>> & {}^0P \\


@VV{\lambda}V & & @AA{UD(\rho)}A \\


\Fr(X)& @>{\f''}>> & {}^0 \Fr(X)


\end{CD}


\qquad \qquad


\begin{CD}


{}^0 P & @>{\p'}>> & P \\


@VV{UD(\lambda)}V & & @AA{\rho}A \\


{}^0 \Fr(X) & @>{\p''}>> & \Fr(X)


\end{CD}


$$


\end{lem}





\begin{pf}


Пусть $X=\{x_j|j\in{\frak J}\}$. Заметим, прежде всего, что


квазиинъективность для гомоморфизма вида


$\mu : \Fr(X)\to {}^0 \Fr(X)$


означает, что композиция $\mu$ с естественным изоморфизмом


$\alpha= \alpha(X) $ есть гомоморфизм вида


$\Fr(X)\to \Fr(X) $, $ x_j\longmapsto x_j+q_j$,


$q_j\in \Fr^{\bullet}(X)$, $j\in{\frak J}$.


Аналогично обстоят дела с квазиинъективными гомоморфизмами из


${}^0 \Fr(X)$ в $\Fr(X)$.


Положим $p_j=\rho(x_j)$,


$j\in{\frak J} $, и пусть


$P^1=\lambda^{-1}(\Fr^1(X))$,


${}^0 P^1={}^{\ast}\lambda^{-1}(\Fr^1(X))$.


Тогда


$P=P^1\oplus P^{\bullet}$,


${}^0 P= {}^0 P^1\oplus {}^0 P^{\bullet}$,


и если ${}^0 p_j={}^{\ast}\rho(x_j) $, то


${}^0 p_j \in {}^0 P^1$. Так как


$P^1\cong D(P)$, ${}^0 P^1 \cong D({}^0 P)$,


то существует запись


$\f'(p_j)={}^0 p_j+h_j({}^0 p)$,


$\p'({}^0 p_j)=p_j+ g_j(p)$,


где $h_j(X)\in \Fr^{\bullet}(X)$,


$h_j({}^0 p)\in {}^0 P^{\bullet}$ и аналогично для $g_j$.





Положим


${\widetilde h}_j=\alpha(UD(\lambda)(h_j({}^0 p)))$,


${\widetilde g}_j=\lambda(g_j(p))$,


$\f''(x_j)=\alpha^{-1}(x_j+{\widetilde h}_j)$,


$\p''(\alpha^{-1} (x_j))=x_j+{\widetilde g}_j$,


и проверим, например, коммутативность левого квадрата. Пусть


$f_j=\lambda(\rho(x_j))$, $j\in{\frak J} $. Тогда


\begin{gather*}


UD(\rho)(\f''(\lambda(p_j))) =


UD(\rho)(\f''(\lambda(\rho(x_j))))= UD(\rho)(\f''(f_j(X)))= \\


%


= UD(\rho)(f_j(\alpha^{-1}(X+\widetilde h)))=


{}^{\ast}\rho(f_j(X+\widetilde h))= f_j({}^0 p + h_j({}^0 p)) = \\


%


= f_j(\f'(p))=\f'(f_j(p))=\f'(p_j).


\end{gather*}


Коммутативность правого квадрата проверяется аналогично,


и, таким образом, лемма доказана.


\end{pf}





\begin{thm}


Пусть $P$ --- проективная алгебра из многообразия $\M$.


\begin{itemize}


\item[(1)] Пусть дан квазиинъективный гомоморфизм


$\f' : P\to {}^0 P$. Тогда найдутся такие


$\pi:\Fr(Y)\to P$,


$\te : P \to \Fr(Y)$, $\pi\cdot\te=1_P $, что


$\f'=\f(\pi,\te)$.


Аналогично, если дан квазиинъективный


$\p' : {}^0 P\to P$, то найдутся


такие $\pi,\te$, что


$\p'=\p(\pi,\te)$.


\item[(2)] Пусть даны два квазиинъективных гомоморфизма


$\p': {}^0 P\to P$, $\f' : P\to {}^0 P$,


Тогда существуют гомоморфизмы


$\pi : \Fr(X')\to P$,


$\te_1,\te_2 : P\to \Fr(X')$


такие, что $\pi\cdot\te_i=1_P$, $i=1,2 $, и


$\f'=\f(\pi,\te_1)$, $\p'=\p(\pi,\te_2)$.


\item[(3)] При тех же предположениях существуют гомоморфизмы


$\pi_1,\pi_2 : \Fr(X'')\to P$,


$\te : P\to \Fr(X'')$


такие, что $\pi_i\cdot\te=1_P$, $i=1,2 $, и


$\f'=\f(\pi_1,\te)$, $\p'=\p(\pi_2,\te)$.


\end{itemize}


\end{thm}





\begin{pf} Пусть вначале


$P=\Fr(X)$, $X=\{x_j|j\in {\frak J}\}$.


Как уже было отмечено выше,


$\alpha(X)(\f'(x_j))=x_j+m_j(X)$,


$\p'(\alpha^{-1}(X)(x_j))=x_j+n_j(X)$,


$m_j,n_j\in \Fr^{\bullet}(X)$, $j\in{\frak J}$.





Докажем сначала п.\,(3). Положим


$Y=X^0 \bigcup X^1 \bigcup X^2 \bigcup X^3 $,


где $X^k$, $k=1,2,3$ ---


непересекающиеся множества одинаковой мощности,


индексированные элементами ${\frak J}$.


Определим


$\te : P\to \Fr(Y)$ так:


$\te(x_j)=x^0_j-x^1_j+m_j(X^2)-n_j(X^3)$, $j\in{\frak J}$.


Здесь $x^k_j\in X^k$. Построим


$\pi_1,\pi_2:\Fr(Y)\to P$, полагая


$\pi_1(x^0_j)=x_j$, $\pi_1(x^1_j)=m_j$,


$\pi_1(x^2_j)=x_j$, $\pi_1(x^3_j)=0$,


$\pi_2(x^0_j)=x_j$, $ \pi_2(x^1_j)=-n_j$,


$\pi_2(x^2_j)=0$, $\pi_2(x^0_j)=x_j$.





Очевидно,что $\pi_1\cdot\te=\pi_2\cdot\te=1_P$. Проверим,что


$\f(\pi_1,\te)=\f'$, $\p(\pi_2,\te)=\p'$.


Из определения естественного изоморфизма $\alpha$


следует, что


$\alpha(X)({}^{\ast} \pi_1(x^0_j))=\alpha(X)({}^{\ast} \pi_1(x^2_j))


=x_j$,


$\alpha(X)({}^{\ast} \pi_1(x^1_j))=\alpha(X)({}^{\ast} \pi_1(x^3_j))


= 0$.


Тогда


$\alpha(X)(\f(\pi_1,\te)(x_j))= \alpha(X)({}^{\ast} \pi_1(\te(x_j)))=


\alpha(X)({}^{\ast} \pi_1 (x^0_j-x^1_j+m_j(X^2)-n_j(X^3)))=


x_j+m_j(X) =\alpha(X)(f'(x_j))$.


Аналогично делается проверка для $\p'$.





Если $P$ --- произвольная проективная алгебра, то


выберем для нее произвольную ретракцию


$\rho : \Fr(X)\to P$, $\lambda : P


\to \Fr(X)$, $\rho\cdot \lambda=1_P$,


и применим лемму 3, чтобы получить квазиинъективные


$\f''$, $\p''$. К этим гомоморфизмам применим


только что доказанный факт:


$\f''=\f(\pi_1,\te),\quad \p''=\p(\pi_2,\te)$.


Теперь по лемме 2


$\f'=\f(\rho\cdot\pi_1,\te\cdot\lambda)$,


$\p'=\p(\rho\cdot\pi_2,\te\cdot\lambda)$.





Пункт (2) доказывается примерно по такому же плану. Вначале


$P=\Fr(X)$, $X=\{x_j|j\in {\frak J}\}$.


$\alpha(X)(\f'(x_j))=x_j+m_j(X)$,


$\p'(\alpha^{-1}(X)(x_j))=x_j+n_j(X)$,


$m_j,n_j\in \Fr^{\bullet}(X)$, $j\in{\frak J}$.


Введем три непересекающиеся множества,равномощные $X$ :


$X^0$, $X^1$, $X^2$, и пусть


$Y=X^0\bigcup X^1 \bigcup X^2 $,


$\pi: \Fr(Y)\to \Fr(X)$


строится так:


$\pi(x^0_j)= x_j-m_j(X)$,


$\pi(x^1_j)= m_j(X)$,


$\pi(x^2_j)= n_j(X)$.


Определим расщепляющие $\pi$ мономорфизмы


$\te_1, \te_2 : \Fr(X) \to \Fr(Y)$ следующим образом:


$$


\te_1(x_j)=x^0_j+m_j(X^0+X^1),\quad


\te_2(x_j)=x^0_j+x^1_j+x^2_j-n_j(X^0+X^1).


$$


Здесь $X^0+X^1$, разумеется, означает,


что вместо каждого $x_j$ подставлены $x^0_j+x^1_j$.


Легко убедиться,что


$\pi\cdot\te_1 =\pi\cdot\te_2=1_P$.


Кроме того,


$\alpha(X)({}^{\ast}\pi(x^0_j))=x_j$,


$\alpha(X)({}^{\ast}\pi(x^1_j))= \alpha(X)({}^{\ast}\pi(x^2_j))= 0$,


${}^{\ast}\te(\alpha^{-1}(X)(x_j))=x^0_j+x^1_j+x^2_j$,


откуда следует, что


$\f(\pi,\te_1)={}^*\pi\cdot\te_1=\f'$,


$\p(\pi,\te_2)=\pi\cdot{}^*\te_1=\p'$.





Для произвольного $P$ поступаем


точно так же, как и при доказательстве п.\,(3).





Доказательство п.\,(1) теперь почти очевидно. Достаточно


например, рассмотреть $P=\Fr(X)$,


$Y= X^0\bigcup X^1\bigcup X^2$,


$\pi(x^0_j)=x_j $, $\pi(x^1_j)=-m_j$,


$\pi(x^2_j)=x_j$,


$\te(x_j)=x^0_j+x ^1_j+m_j(X^2).$ Тогда $\f'=\f(\pi,\te)$.


\end{pf}





Поскольку $\Fr(X)\cong UD\Fr(X) $, вопрос о том, каждая ли


проективная алгебра изоморфна алгебре вида $\Fr(X)$, естественно


обобщить так: каждая ли алгебра $ P$ изоморфна алгебре вида


$UD(Q)$. Если $ P\cong UD(Q) $, то ввиду того,


что $DU(N)\cong N $, имеем $D(Q)\cong D(P)$,


и, таким образом


$P\cong UD(Q)\cong UD(P)={}^0 P$.


Следовательно, ${}^0 P$ можно считать алгеброй,


наиболее близкой и к $P$, и к свободной алгебре.





Будем называть гомоморфизм


$\nu :UD(B)\to UD(C)$ однородным, если


$\nu=U(\beta(D(C))^{-1})\cdot UD(\nu) \cdot U(\beta(D(B)))$.


В частности, однородны все гомоморфизмы вида $U(\gamma)$.


Это следует из естественности $\beta$.





\begin{cor}


Пусть имеется гомоморфизм проективных алгебр


$\mu : P\to {}^0 P$


такой, что $D(\mu)$ --- изоморфизм. Тогда существует однородный


автоморфизм $\lambda : {}^0 P\to {}^0 P $,


и гомоморфизм $\f(\pi,\te) : P\to {}^0 P$


такие, что


$\mu=\lambda\cdot \f(\pi,\te)$.


Аналогичное утверждение имеет место для гомоморфизмов из ${}^0 P$


в $P$ с заменой $\f$ на $\p$.


\end{cor}





\begin{pf} Положим


$\lambda=U(\beta(D(P)))^{-1}\cdot UD(\mu)$.


Ясно, что это --- автоморфизм.


Проверим однородность $ \lambda$, которая эквивалентна равенству


$\lambda=U(\beta(D(P)))^{-1}\cdot UD(\lambda)\cdot U(\beta(D(P)))$.


Оно будет следовать из такого факта:


$$


D(\mu)=DU(\beta(D(P))^{-1})\cdot DUD(\mu)\cdot\beta(D(P)).


$$





Доказательство этого соотношения сводится к двум тождествам


--- следствиям естественности $\beta $: 


$DUD(\mu)\cdot\beta(D(P))=\beta(DUD(P))\cdot D(\mu)$,


$DU(\beta(D(P)))\cdot \beta(D(P))=\beta(DUD(P))\cdot \beta(D(P))$.





В частности, $DU(\beta(D(P)))= \beta(DUD(P))$.


Теперь рассмотрим гомоморфизм $\lambda^{-1}\mu$


и вычислим $ D(\lambda^{-1}\mu)$,


используя полученное выше выражение для $D(\mu)$:


\begin{gather*}


D(\lambda^{-1}\mu)=DUD(\mu)^{-1}\cdot DU(\beta(D(P)))\cdot D(\mu)=\\


=DUD(\mu)^{-1}\cdot\beta(DUD(P))\cdot D(\mu)= \\


=DUD(\mu)^{-1}\cdot DUD(\mu)\cdot \beta(D(P))=\beta(D(P)).


\end{gather*}





Остается применить теорему 2: ${\lambda}^{-1}\mu=\f(\pi,\te)$.


\end{pf}





Далее нам потребуются естественные изоморфизмы:


$\alpha_1 : U\Fr_{\R} \ar{\cong} \Fr_{\M}$,


$\alpha_2 : D\Fr_{\M} \ar{\cong} \Fr_{\R}$


такие, что $\alpha =\alpha_1\cdot U(\alpha_2)$,


и одно соотношение между ними: композиция


$$


D\Fr(X)@>{\beta(D\Fr(X))}>>


DUD\Fr(X) @>{D(\alpha(X))}>> D\Fr(X)


$$


есть тождественный изоморфизм. В интересующей нас ситуации все это


имеется.





Определим функтор на категории


свободных $\T$--алгебр, который тождественен на объектах,


а гомоморфизмы $\nu: \Fr(X)\to \Fr(Y)$ отображает в


$$


\o{\nu}: \Fr(X)\ar{\alpha^{-1}(X)}UD\Fr(X)\ar{UD(\nu)}


UD\Fr(Y)\ar{\alpha(Y)} \Fr(Y).


$$


Легко убедиться, что


$\o{\o{\nu}}=\o{\nu}$.


Очевидно также, что $\o{\te\cdot\pi}=\a{\te}\cdot\a{\pi}$,


и что


$\f(\pi,\te)=\a{\pi}\cdot\te$, $\p(\pi,\te)=\pi\cdot(\a{\te})$.





Сформулируем условия изоморфизма $\f(\pi,\te)$ и


$\p(\pi,\te)$.


Пусть даны $e,\ep:\Fr(X)\to \Fr(X)$,


причем $e^2=e $.


$$


\begin{array}{lccc}


&\text{Условие } {\text{ISO}}(e,\ep,\f) & \hspace{1cm} &


\text{Условие } {\text{ISO}}(e,\ep,\p) \\


{\text{ISO\,1}} & \o{e}e\ep=\o{e} & & \ep e\o{e}=\o{e}\\


{\text{ISO\,2}} & e\ep e=e & & e\ep e=e \\


{\text{ISO\,3}} & \ep\o{e}=\ep & & \o{e}\ep=\ep


%


% \text{Условие }&{\text{ISO}}(e,\ep,\f) & \hspace{1cm}


% \text{Условие }&{\text{ISO}}(e,\ep,\p) \\


% {\text{ISO\,1}} \o{e}e\ep&=\o{e} & \ep e\o{e}&=\o{e}\\


% {\text{ISO\,2}} e\ep e&=e & e\ep e&=e \\


% {\text{ISO\,3}} \ep\o{e}&=\ep & \o{e}\ep&=\ep


\end{array}


$$





\begin{thm} Пусть дана ретракция


$\pi: \Fr(X)\to P$, $\te : P\to \Fr(X)$,


$\pi\cdot\te=1_P$, $ e=\te\cdot\pi$.


Эквивалентны следующие условия:


\begin{itemize}


\item[(1)] $\f(\pi,\te)$ --- изоморфизм,


\item[(2)] выполнено условие ${\mathrm{ISO}}(e,\ep,\f) $


для некоторого $ \ep $,


\item[(3)] существует гомоморфизм $u:\Fr(X)\to \Fr(X)$,


для которого


$u= e\cdot u=u\cdot\o{e}$, $u\cdot e=e$, $\o{e}\cdot u= \o{e}$.


\end{itemize}





Гомоморфизм $u$ в п.\,{\em (3)} определяется приведенными


там условиями однозначно, причем выполняются соотношения


$u=u^2$, $\o{u}=\o{e}$. При этом


$\f(\pi,\te)^{-1}=\pi\cdot u\cdot(\a{\te})$.


По заданному $\ep $ из п.\,{\em (2) } также можно найти


$\f(\pi,\te)^{-1}=\pi\cdot \ep\cdot(\a{\te})$,


а связь между $u$ и $\ep$ такова: $u=e\cdot\ep$.





Условия обратимости $\p(\pi,\te) $ формулируются двойственным образом.


\end{thm}





\begin{pf} (1)\,$\Longrightarrow$\,(3). Пусть


существует $h=\f^{-1}$, $\f=\f(\pi,\te)$. Положим


$u=\te\cdot h\cdot(\a{\pi})$, тогда $h=\pi\cdot u\cdot


(\a{\te})$. Проверим условия п.\,(3):


\begin{gather*}


u\cdot e=\te\cdot h\cdot\a{\pi}\cdot\te\cdot\pi=


\te\cdot\f^{-1}\cdot \f\cdot\pi=\te\cdot\pi=e,\\


\o{e}\cdot u=\a{\te}\cdot\a{\pi}\cdot\te\cdot h \cdot\a{\pi}=


\a{\te}\cdot\f\cdot\f^{-1}\cdot\a{\pi}=\o{e}.


\end{gather*}


Остальное еще более тривиально.


Покажем единственность. Пусть существуют гомоморфизмы


$u_1$ и $u_2$, удовлетворяющие соотношениям п.\,(3).


Тогда


$u_1 u_2=(u_1 \o{e})u_2=u_1(\o{e}u_2)=u_1 \o{e}=u_1$,


$u_1 u_2=u_1(e u_2)=(u_1 e)u_2=eu_2=u_2$.





(3)\,$\Longrightarrow$\,(2). Из условий п.\,(3) следует


выполнимость условия ${\mathrm{ISO}}(e,u,\f):$


$\o{e}(eu)=\o{e}u=\o{e}$, $ e(ue)=ee=e$, $u\o{e}=u$.





(2)\,$\Longrightarrow$\,(1). Положим $h=\pi\cdot\ep(\a{\te})$


и проверим, что $\f\cdot h=1$, $h\cdot \f=1$.


$$


\f\cdot h=1 \Longleftrightarrow


\a{\te}\cdot(\f\cdot h)\cdot \a{\pi}=\a{\te}\cdot\a{\pi}


\Longleftrightarrow \o{e}\cdot e\cdot \ep\cdot \o{e}=\o{e},


$$


но последнее соотношение эквивалентно условию ${\mathrm{ISO\,1}}$.


Далее,


$h\cdot\f=\pi\cdot\ep\cdot\o{e}\cdot\te=\pi\cdot\ep\cdot\te$,


т.\,к. по ${\mathrm{ISO\,3}}\;$ $\ep\cdot\o{e}=\ep$. Теперь


$$


h\cdot\f=1 \Longleftrightarrow


\te\cdot(h\cdot\f)\cdot\pi=\te\cdot\pi


\Longleftrightarrow


e\cdot\ep\cdot e= e,


$$


т.\,е. все сводится к условию ${\mathrm{ISO\,2}}$.





Наконец, пусть выполнено условие ${\mathrm{ISO}}(e,\ep,\f)$. Положим


$u=e\cdot\ep$ и проверим выполнимость соотношений п.\,(3):


$eu=ee\ep=e\ep=u$, $u\o{e}=e(\ep\o{e})=e\ep=u$,


$ue=e\ep e=e$, $\o{e} u=\o{e}e\ep=\o{e}$.





Утверждения, относящиеся к обратимости $\p(\pi,\te)$,


доказываются совершенно аналогично.


\end{pf}





\begin{cor} Если многообразие $\Alg({\T})$ таково, что


все $\f(\pi,\te)$ инъективны (а такими являются, например,


все однородные многообразия (см. \cite{T2}, \cite{T4}), то


для обратимости $\f(\pi,\te)$ необходимо и достаточно


выполнение условия ${\mathrm{ISO\,1}}$ из набора условий


${\mathrm{ISO}}(e,\ep,\f)$. Аналогичное утверждение имеет место


и для $\p({\pi,\te})$.


\end{cor}





Доказательство по сути дела содержится в части


(2)\,$\Longrightarrow$\,(1) доказательства теоремы 1.





\begin{cor} Для того чтобы $\p(\pi,\te)=\f(\pi,\te)^{-1}$,


необходимо и достаточно, чтобы выполнялись соотношения


$\eo e \eo = \eo $, $e \eo e = e$.


Если многообразие однородно, то эти соотношения равносильны.


\end{cor}





Это легко следует из теоремы 3, но еще легче доказывается непосредственно.





\begin{cor}


Пусть дана ретракция


$ \pi : \Fr(X)\to P$, $ \te : P\to \Fr(X)$,


$\pi\cdot\te=1_P $,$ e=\te\cdot\pi$.


Если выполнено условие ${\mathrm{ISO}}(e,\ep,\f) $ при


$\ep=\o{\ep}=\ep^2$, то выполнены условия предыдущего


следствия, и кроме того существуют обратимые слева вложения


вида


\begin{align*}


% \la' &: UD(P)\ar{UD(\te')}UD\Fr(X)\ar{\alpha(X)} \Fr(X),\\


% \la'' &: UD(Q)\ar{UD(\te'')}UD\Fr(X)\ar{\alpha(X)} \Fr(X),


\la' &: UD(P) @>{UD(\te')}>> UD\Fr(X) @>{\alpha(X)}>> \Fr(X),\\


\la'' &: UD(Q) @>{UD(\te'')}>> UD\Fr(X) @>{\alpha(X)}>> \Fr(X),


\end{align*}


индуцирующие изоморфизм


$UD(P) \coprod UD(Q)\ar{\cong} \Fr(X)$,


и такие, что композиция $\pi\cdot\la'$ есть гомоморфизм,


обратный к $\f(\pi,\te)$.


\end{cor}





Первое утверждение следствия 4 очевидно.


Прежде чем доказывать все остальное,


сделаем несколько более общих наблюдений.


Назовем $\nu : \Fr(X) \to \Fr(Y) $ однородным, если $\o{\nu}=\nu$


(``однородность'', введенная перед следствием 1 --- почти то же самое,


к тому же она больше не будет употребляться).





\begin{lem}


Гомоморфизм $\nu : \Fr(X)\to \Fr(Y)$ однороден тогда и


только тогда, когда


$\nu=\alpha_1(Y)\cdot U(\la)\cdot {\alpha_1}^{-1}(X)$.


\end{lem}





\begin{pf} Если $\nu$ однороден, то


$$


\mu=\o{\nu}=\alpha UD(\nu)\alpha^{-1}=


\alpha_1 U(\alpha_2)UD(\nu)U(\alpha_2^{-1})\alpha_1^{-1}=


\alpha_1 U(\alpha_2 D(\nu)\alpha_2^{-1})\alpha_1^{-1}.


$$


Обратно, пусть $\nu=\alpha_1(Y)\cdot U(\la)\cdot \alpha_1^{-1}(X)$.


Заметим сначала, что условие $D(\alpha)\cdot\beta(D\Fr)=1$


сводится к


$D(\alpha_1(X))\beta(\Fr_{\R}(X))\alpha_2(X)=1$,


откуда вытекает


$\beta(\Fr_{\R}(X))=D(\alpha_1^{-1}(X))\alpha_2^{-1}(X)$.\linebreak


Теперь коммутативность диаграммы, соответствующая равенству


$\nu=\o{\nu}$, следует из естественности~$\beta$.


\end{pf}





Из этой леммы следует, что однородные гомоморфизмы свободных алгебр ---


это в точности гомоморфизмы, отображающие базисные элементы


$x_j$ в элементы вида $\sum\limits_i y_ia_{ij}$, $a_{ij}\in R$.


В частности, таковы гомоморфизмы вида $ \Fr(\gamma) $, например,


гомоморфизм $\In_X:\Fr(X)\to \Fr(X\cup Y)$, соответствующий


вложению $X\subseteq X\cup Y$.





\begin{lem}


Из условия ${\mathrm{ISO}}(e,\ep,\f)$, $e=\te\cdot\pi$,


$\ep=\ep^2=\o{\ep}$


следует,что


$\ep=\te'\cdot\a{\pi}$,


$\a{\pi}\cdot\te'=1$, $\te'=\a{(\ep\cdot\te)}$.


Двойственным образом, из ${\mathrm{ISO}}(e,\ep,\p) $ и


$\ep=\ep^2=\o{\ep}$


следует,


что


$\ep=\a{\te}\cdot\pi'$, $\pi'\cdot\a{\te}=1$,


$\pi'=\a{(\pi'\cdot\ep)}$.


\end{lem}





\begin{pf} Пусть $\epo=\ep=\ep^2$. Положим


$\te'=\ep\cdot\a{\te}$.


Тогда в силу однородности $\ep$


$$


\te'=\ep\cdot\alpha(X)\cdot UD(\te)=


\alpha(X)\cdot UD(\ep)\cdot UD(\te)=


\alpha(X)\cdot UD(\ep\cdot\te)= \a{(\ep\te)}.


$$


Рассмотрим $\a{\pi}\cdot \te'=\a{\pi}\cdot\ep\cdot\a{\te}$.


Чтобы показать, что это тождественный гомоморфизм, умножим его


слева на мономорфизм $\a{\te}$ и справа на эпиморфизм


$\a{\pi}$. Получится $\o{e}\ep\o{e}$, что по


условию ${\mathrm{ISO\,2}}$ равно


$\o{e}=\a{\te}\cdot\a{\pi}$. Сокращая на мономорфизм


слева и эпиморфизм справа, получаем требуемое.


\end{pf}





Теперь можно закончить доказательство следствия 4.


Обратный к $\f(\pi,\te)$ гомоморфизм имеет вид


$\pi\cdot\ep\cdot(\a{\te})=\pi\cdot(\a{\ep\cdot\te})$.


Положим $\te'=\ep\cdot\te$, тогда


$\ep = \a{\te'}\cdot (\a{\pi})$, $\a{\pi}\cdot (\a{\te'})=1$.


Ввиду однородности $\ep$ этот гомоморфизм отображает


элементы базиса свободной алгебры $x_j$ в элементы


вида $\sum\limits_i x_i a_{ij} $, причем матрица $E$


с компонентами $a_{ij}$ идемпотентна. Рассмотрим


идемпотентный эндоморфизм $\ep'$ свободной алгебры $\Fr(X)$,


отображающий элементы базиса $x_j$ в элементы


$\sum\limits_i x_i(\delta_{ij}-a_{ij})$, где $\delta_{ij}$


--- символ Кронекера. Образ $Q$ этого эндоморфизма удовлетворяет


всем необходимым условиям. Следствие доказано.





В связи с только что доказанным утверждением отметим несколько


работ, касающихся близких вопросов. В \cite{Bur} доказано, что


в шрейеровых многообразиях над полями для каждой сюръекции свободных алгебр


$\pi:\Fr(X\cup Y)\to \Fr(X)$


существует автоморфизм $\gamma$ алгебры $\Fr(X\cup Y)$


такой, что композиция


$$


\Fr(X)\ar{\In_X} \Fr(X\cup Y)\ar{\gamma} \Fr(X\cup Y)\ar{\pi} \Fr(X)


$$


есть тождественный гомоморфизм, а композиция


$$


\Fr(Y)\ar{\In_Y} \Fr(X\cup Y)\ar{\gamma} \Fr(X\cup Y)\ar{\pi} \Fr(X)


$$


есть нулевой гомоморфизм.


Там же дан пример сюръекции


колец многочленов над полем $R$ характеристики 2


$\pi: R[x_1,x_2,x_3]\to R[y_1,y_2]$


со следующим свойством: не существует


$z_1,z_2,z_3\in R[x_1,x_2,x_3]$ таких, что


$$


R[z_1,z_2,z_3]= R[x_1,x_2,x_3],\quad


\pi(z_1)=y_1,\quad \pi(z_2)=y_2,\quad \pi(z_3)=0.


$$


В статье \cite{A4} построена ретракция свободной метабелевой алгебры


Ли, образ которой не выделяется свободным множителем. Напомним, что


в \cite{A2} было показано что любой ретракт свободной метабелевой


алгебры Ли над полем изоморфен свободной алгебре.





\begin{thm}


Пусть дан квазиинъективный $\f': P\to {}^0 P$.


Этот гомоморфизм является изоморфизмом


тогда и только тогда, если существует такая ретракция


$\pi_{+}: \Fr(Z)\to P$, $\te_{+}: P\to \Fr(Z)$,


$\pi_{+}\cdot\te_{+}=1_P$,


что $\f'=\f(\pi_{+},\te_{+})$, и для


$e_{+}=\te_{+}\cdot\pi_{+}$ выполнено условие


${\mathrm{ISO}}(e_{+},\ep_{+},\f)$, причем {\em (}и это существенно{\em !)}


$\o{\ep_{+}}=\ep_{+}=\ep^2_{+} $. В частности, выполнены условия


следствия {\em 4}. Аналогичный факт имеет место для квазиинъективного


$\p': {}^0 P\to P$.


\end{thm}





Утверждение для $\p'$ следует из соответствующего факта для


$\f'$, т.к. $(\p')^{-1}=\f'$ --- квазиинъективен, причем по


следствию 3 $\;\f'=\f(\pi_{+},\te_{+})$, $\p'=\p(\pi_{+},\te_{+})$


и выполнены условия $\o{e_{+}}, e_{+},\o{e_{+}}=\o{e_{+}}$,


$e_{+},\o{e_{+}}, e_{+}=e_{+}$,


т.\,е. выполнено $ {\mathrm{ISO}}(e_{+},\ep_{+},\p)$.


Следовательно,достаточно разобрать случай $\f'$.


Достаточность условий теоремы 4 следует из теоремы 3, а необходимость


будет следовать из следующего более конкретного факта.





\begin{thm} Пусть даны $\pi : \Fr(X)\to P$, $ \te : P\to \Fr(X)$,


$\pi\cdot \te =1_P$.


$\;\f(\pi,\te)$ --- изоморфизм тогда и только тогда, когда


существуют непересекающееся с $X$ множество $Y$,


возможно, даже пустое


{\em (}конечное, если $X$ конечно{\em )}, и ретракция


$\pi_{+} : \Fr(X\cup Y)\to P$,


$\te_{+}: P\to \Fr(X\cup Y)$, $\pi_{+}\cdot\te_{+}=1_P$


такие, что если $\In_X: \Fr(X)\to \Fr(X\cup Y)$ --- естественный


гомоморфизм, соответствующий включению $X\subseteq X\cup Y$,


и $e_{+}=\te_{+}\cdot\pi_{+}$, то


$\pi_{+}\cdot \In_X=\pi $, $\In_X\cdot\te=\te_{+}$,


и для некоторого $\ep_{+}=\ep_{+}^2$ выполнено условие


${\mathrm{ISO}}(e_{+},\ep_{+},\f)$, причем


$\o{\ep_{+}}=\ep_{+}$.


\end{thm}





{\bf Доказательство} теоремы будет представлено в виде цепочки лемм.


Легко проверяется справедливость следующего утверждения.





\begin{lem}


Путь дана коммутативная диаграмма


$$


\begin{CD}


\Fr(X) @>{\pi_1}>> P @>{\te_1}>> \Fr(X) \\


@VV{\nu}V 


   \begin{picture}(0,20)


   \put(-40,-5){\vector(2,1){35}}


   \put(-35,8){$\pi_2$}


   \put(5,13){\vector(2,-1){35}}


   \put(25,8){$\te_2$}


   \end{picture}


&& @VV{\nu}V \\


\Fr(Y) & & & & \Fr(Y)


\end{CD}


$$


где $\pi_i\te_i=1_P$, $i=1,2$ и $\nu$ однороден.


Тогда


$\f(\pi_1,\te_1)= \f(\pi_2,\te_2)$, $\p(\pi_1,\te_1)= \p(\pi_2,\te_2)$.


\end{lem}





Из этой леммы и из теоремы 3 следует достаточность в теореме 5.


Каждую из следующих лемм можно сформулировать и доказать


для двух двойственных случаев.


Формулировки и доказательства даны только для случая $\f$,


т.к. в случае $\p$ выкладки видоизменяются очевидным образом.





\begin{lem}


Пусть даны две ретракции


$$


\rho_j:\Fr(X_j)\to Q, \quad \la_j: Q\to \Fr(X_j),\quad


\rho_j\la_j=1_Q,\quad j=1,2.


$$


Положим $u_j=\la_j\cdot\rho_j$,


$p=\la_2\rho_1 $, $ q=\la_1\rho_2$.


Тогда равносильны следующие условия:


\begin{itemize}


\item[(1)] $\f(\rho_1,\la_1)=\f(\rho_2,\la_2)$;


\item[(2)] $\o{u_1}\cdot u_1=\o{q}\cdot p$;


\item[(3)] $\o{u_2}\cdot u_2=\o{p}\cdot q$.


\end{itemize}


\end{lem}





\begin{pf} Из условий следует, что


$u_1=q\cdot p$, $u_2=p\cdot q$. Докажем равносильность


пп.\,(1) и (2). Пункт (1) фактически есть равенство


$\a{\rho_1}\cdot\la_1=\a{\rho_2}\cdot\la_2$


Умножая его слева на мономорфизм $\a{\la_1}$,


а справа на эпиморфизм $\rho_1$, получим равенство из п.\,(2).


Очевидно, что эти действия обратимы, т.к. на мономорфизм


слева и на эпиморфизм справа можно сокращать. Точно так же


устанавливается равносильность (1) и (3).


\end{pf}





\begin{lem}


Пусть дана коммутативная диаграмма


$$


\begin{CD}


\Fr(X_1) @>{\rho_1}>> Q @>{\la_1}>> \Fr(X_1) \\


@VV{\nu}V 


   \begin{picture}(0,20)


   \put(-40,-5){\vector(2,1){35}}


   \put(-35,8){$\rho_2$}


   \put(5,13){\vector(2,-1){35}}


   \put(25,8){$\la_2$}


   \end{picture}


&& @VV{\nu}V \\


\Fr(X_2) & & & & \Fr(X_2)


\end{CD}


$$


в которой $\rho_i\la_i=1_Q$, $i=1,2$, и $\nu$


однороден. Положим $u_i= \la_i\cdot\rho_i$, $i=1,2$. Тогда,


если выполнено условие ${\mathrm{ISO}}(u_1,\ep',\f)$ для


$\ep'=\o{\ep'}=\ep^2$, то выполнено условие


${\mathrm{ISO}}(u_2,\ep'',\f)$ при


$\ep''=\nu\cdot \ep'\cdot \o{q}$, где $q=\la_1\rho_2$,


причем $\ep''=(\ep'')^2=\o{\ep''}$.


\end{lem}





{\bf Доказательство.} Положим $p=\la_2\rho_1$,


Тогда имеют место следующие соотношения:


\begin{gather*}


u_1=qp=q\nu \quad u_2=pq=\nu q \quad qpq=q\quad pqp=p, \\


\o{u_1} u_1=\o{q} p\quad \o{u_2} u_2=\o{p} q, \\


p= u_2\nu=\nu u_1 \quad


q=q u_2=u_1 q.


\end{gather*}


Рассмотрим $\ep''=\nu\cdot \ep'\cdot \o{q}$.


Очевидно, что $\ep''=\o{\ep''}$. Идемпотентность следует из


равенств, при выводе которых применяются предыдущие соотношения и


условия ${\mathrm{ISO}}(u_1,\ep',\f)$:


$$


(\ep'')^2=\nu\cdot \ep'\cdot (\o{q}\cdot\nu)\cdot\ep'\cdot \o{q} =


\nu\cdot \ep'\cdot (\o{u_1}\cdot\ep')\cdot \o{q}


= \nu\cdot \ep'\cdot \o{u_1}\cdot \o{q} =


\nu\cdot \ep'\cdot \o{q} = \ep''.


$$





Условие ${\mathrm{ISO\,3}}$ выводится так


$$


\ep''\cdot \o{u_2}=\nu\cdot\ep'\cdot\o{q}\cdot\o{u_2}=


\nu\cdot\ep'\cdot\o{u_1}\cdot\o{q} = \nu\cdot\ep'\cdot\o{q}.


$$





Выведем условие ${\mathrm{ISO\,1}}$


$$


\o{u_2} u_2 \ep''= \o{ u_2} u_2 \nu \ep' \o{q}=


\o{p} q \nu \ep' \o{q}=


\o{p} u_1 \ep'\o{q}=\nu \o{u_1} u_1 \ep' \o{q}=


\nu\cdot\o{u_1}\cdot \o{q}=\o{p}\cdot\o{q}=\o{u_2}.


$$





Наконец, условие ${\mathrm{ISO}\,2} $ получается так


$$


% break into miltline* in Eng.


u_2\ep''u_2 = u_2\nu\ep'\o{q}u_2= \nu u_1\ep'\o{q}(\o{u_2}u_2)=


\nu u_1\ep'\o{q}(\o{p}q)=\nu u_1\ep'\o{q}\o{u_1}q=


\nu u_1\ep'q=\nu u_1\ep' u_1 q=\nu u_1 q=pq=u_2.\quad\square


$$





\begin{lem}


Пусть даны гомоморфизмы


$$


\rho: \Fr(Z)\to Q,\quad \gamma_1,\gamma_2: Q\to \Fr(Z),\quad


\rho\cdot\gamma_i=1_Q,\quad i=1,2.


$$


Допустим, что $\f(\rho,\gamma_1)=\f(\rho,\gamma_2)$.


Положим $u_i=\rho\cdot\gamma_i$, $i=1,2$.


Тогда, если выполнено условие ${\mathrm{ISO}}(u_1,\ep,\f)$,


причем $ \epo=\ep$, то выполнено и условие ${\mathrm{ISO}}(u_2,\ep,\f)$.


\end{lem}





\begin{pf} Из условий следует, что


$u_1=u_1u_2$, $u_2=u_2 u_1$. В обозначениях леммы 7


$p=u_2$, $q=u_1$ и условие


$\f(\rho_1,\gamma)=\f(\rho_2,\gamma)$ эквивалентно


любому из равенств


$\o{u_1}u_1= \o{u_1}u_2$, $\o{u_2}u_2=\o{u_2}u_1$.


Дальнейшие вычисления аналогичны проделанным выше. Проверим,


например, что $u_2\ep u_2=u_2$. Это равенство эквивалентно


равенству $\rho \ep \gamma_2=1_Q$ (сокращаем на мономорфизм


слева и на эпиморфизм справа). Согласно лемме 5 представим $\ep$


в виде


$\ep=\gamma'\cdot\a{\rho}$, $\a{\rho}\cdot\gamma'=1$,


$\gamma'=\a{(\ep\gamma_1)}$.


Отсюда получаем


$\rho\ep\gamma_2 = \rho\gamma'(\a{\rho})\gamma_2=


(\rho\gamma')\cdot\f(\rho,\gamma_2)=


(\rho\gamma')\cdot\f(\rho,\gamma_1)$.


Но по теореме 3 $\;\rho\gamma'=\f(\rho,\gamma_1)^{-1} $, что и


доказывает наше утверждение. Условия ${\mathrm{ISO\,2}}$ и


${\mathrm{ISO\,3}} $ получаются еще более просто:


$\o{u_2}u_2\ep=\o{u_2}u_1\ep=(\o{u_2}\cdot \o{u_1})$,


$u_1\ep= \o{u_2}(\o{u_1} u_1\ep ) = u_2\cdot\o{u_1}=\o{u_2}$,


$ep\o{u_2}=(\ep\o{u_1})\o{u_2}=\ep(\o{u_1}\o{u_2})=\ep\o{u_1}=\ep$.


\end{pf}





Допустим теперь, что $\f(\pi,\te)$--- изоморфизм. По теореме 3


можно найти множество $Y$ и гомоморфизмы


$\pi' : \Fr(Y)\to P$, $\te_1,\te_2 : P\to \Fr(Y)$,


$\pi'\cdot\te_1=\pi'\cdot\te_2= 1_P$


такие, что


$\f(\pi,\te)=\f(\pi',\te_1)$, $\f(\pi,\te)^{-1}=\p(\pi',\te_2)$.





\begin{lem}


Положим $e_1=\te_1\pi'$, $e_2=\te_2\pi'$. Тогда выполнено


условие ${\mathrm{ISO}}(e_1,\o{e_2},\f)$.


\end{lem}





\begin{pf} По построению


$e_1 e_2=e_1$ и $ e_2 e_1=e_2$. Отсюда


$\o{e_1}\cdot \o{e_2}=\o{e_1}$ и $\o{e_2}\cdot\o{e_1}=\o{e_2}$.


Из взаимной обратности $\f(\pi',\te_1)=\a{\pi'}\cdot\te_1$


и $\p(\pi',\te_2)=\pi'\cdot\a{\te_2}$ получаем


$\a{\pi'}\cdot\te_1\cdot\pi'\cdot\a{\te_2}=1$,


$\pi'\cdot\a{\te_2}\cdot\a{\pi'}\cdot\te_1=1$.


Первое равенство умножаем слева на $\a{\te_1}$,


справа --- на $\a{\pi'}$,


второе --- слева на $\te_1$, справа на $\pi'$.


В результате получим


$\o{e_1}\cdot e_1\cdot\o{e_2}=\o{e_1}$,


$ e_1\cdot\o{e_2}\cdot e_1=e_1$.


Полученные соотношения вместе составляют условие


${\mathrm{ISO}}(e_1,\o{e_2},\f)$.


\end{pf}





Можно считать,что $ X\cap Y=\emptyset $.


Рассмотрим естественные вложения и проекции


\begin{align*}


\In_X&: \Fr(X)\to \Fr(X\cup Y), \quad \pr_X: \Fr(X\cup Y)\to \Fr(X),\\


\In_Y&: \Fr(Y)\to \Fr(X\cup Y),\quad \pr_Y: \Fr(X\cup Y)\to \Fr(Y),


\end{align*}


где, например, $\In_X$ соответствует включению


$ X\subset Y$, а $\pr_X$ отображает базисные элементы


из $X$ в самих себя, а базисные элементы из $Y$ ---


в нуль. Очевидно, что все эти гомоморфизмы однородны в смысле данного


выше определения. Определим $\pi_{+}:\Fr(X\cup Y)\to P$


условиями:


$\pi_{+}\cdot \In_X=\pi$, $\pi_{+}\cdot \In_Y=\pi'$,


и пусть


$\te_{+}=\In_X\cdot \te $, ${\te'}_{+}=\In_Y\cdot \te_1$.


Очевидно, что $\pi_{+}\cdot\te_{+}=\pi_{+}\cdot {\te'}_{+}=1_P$.





Теперь можно закончить доказательство теоремы.


Рассмотрим коммутативную диаграмму


\begin{equation}


\begin{CD}


\Fr(Y) @>{\pi'}>> P @>{\te_1}>> \Fr(Y) \\


@VV{\In_Y}V 


   \begin{picture}(0,20)


   \put(-40,-5){\vector(2,1){35}}


   \put(-35,8){$\pi_{+}$}


   \put(5,13){\vector(2,-1){35}}


   \put(25,8){${\te'}_{+}$}


   \end{picture}


&& @VV{\In_Y}V \\


\Fr(X\cup Y)&&&&\Fr(X\cup Y)


\end{CD}


\end{equation}


Согласно лемме 10, имеет место условие


${\mathrm{ISO}}(\te_1\cdot\pi',\o{\te_2\cdot\pi'},\f)$.


Лемма 8, примененная к диаграмме (2), показывает, что для


некоторого $\ep_{+}=\o{\ep_{+}}=\ep^2_{+}$ выполняется условие


${\mathrm{ISO}}({\te'}_{+}\cdot\pi_{+},\ep_{+},\f)$.


Воспользуемся тем, что имеет место коммутативная диаграмма


\begin{equation}


\begin{CD}


\Fr(X) @>{\pi}>> P @>{\te}>> \Fr(X) \\


@VV{\In_X}V 


   \begin{picture}(0,20)


   \put(-40,-5){\vector(2,1){35}}


   \put(-35,8){$\pi_{+}$}


   \put(5,13){\vector(2,-1){35}}


   \put(25,8){$\te_{+}$}


   \end{picture}


&& @VV{\In_X}V \\


\Fr(X\cup Y)&&&&\Fr(X\cup Y)


\end{CD}


\end{equation}


По лемме 6


$\f(\pi_{+},{\te'}_{+})=\f(\pi',\te_1)=\f(\pi,\te)=\f(\pi_{+},\te_{+})$.


Остается применить лемму 9, согласно которой из


${\mathrm{ISO}}({\te'}_{+}\cdot\pi_{+},\ep_{+},\f)$ следует


${\mathrm{ISO}}(\te_{+}\cdot\pi_{+},\ep_{+},\f)$. Теорема доказана.





\begin{notenonum} Из доказательства следует, что


явный вид $\ep_{+}$ таков:


$\ep_{+}=\In_Y\cdot\o{\te_2\cdot\pi_{+}}$.


\end{notenonum}








\begin{thebibliography}{99}





\bibitem{T1}


Тронин С.Н. {\em О системах образующих проективных


алгебр} // Изв. вузов. Математика. -- 1984. -- \No\,5. -- С.\,63--72.


\bibitem{T2}


Тронин С.Н. {\em Об одной конструкции в теории проективных


алгебр} // Матем. заметки. -- 1984. -- Т.\,35. -- \No\,5. -- C.\,647--652.


\bibitem{T3}


Тронин С.Н. {\em О ретракциях колец многочленов} //


Изв. вузов. Математика. -- 1985. -- \No\,2. -- C.\,84--85.


\bibitem{T4}


Тронин С.Н. {\em О коммутативных ассоциативных проективных


алгебрах ранга} 2 {\em над совершенным полем} // Матем.


заметки. -- 1987. -- Т.\,41. -- \No\,6. -- C.\,776--780.


\bibitem{T5}


Тронин С.Н. {\em О некоторых свойствах алгебраических теорий


многообразий линейных алгебр} II. {\em Универсальное обращение


гомоморфизмов и расслоенные произведения теорий}. -- Казанск.


гос. ун-т. -- Казань, 1988. -- 38 с. -- Деп. в ВИНИТИ 11.08.88,


\No\,6510-В88.


\bibitem{T6}


Тронин С.Н. {\em О ретракциях свободных алгебр и


модулей.} --- Дис.~$\dotsc$ канд. физ.-матем. наук. -- Кишинев,


1989. -- 105 с.


\bibitem{Dn}


{\em Днестровская тетрадь. Нерешенные проблемы теории колец и


модулей}. -- 4-е изд. -- Новосибирск, 1993. -- 73 с.


\bibitem{A1}


Артамонов В.А. {\em Нильпотентность, проективность, свобода}


// Вестн. Моск. ун-та. Сер. матем., механ. -- 1971. -- \No\,5. -- C.\,50--53.


\bibitem{A2}


Артамонов В.А. {\em Проективные метабелевы группы и алгебры


Ли} // Изв. АН СССР. Сер. матем. -- 1978. -- Т.\,42. -- \No\,2. --


C.\,226--236.


\bibitem{A3}


Артамонов В.А. {\em Проективные метабелевы $D$-группы и


супералгебры Ли} // Тр. семин. им. И.Г.\,Петровского. -- 1991. -- 


Вып.\,15. -- С.\,189--195.


\bibitem{Law}


Lawvere F.W. {\em Functorial semantics of algebraic


theories} // Proc. Nat. Acad. Sci. USA. -- 1963. -- V.\,50. -- \No\,5. --


P.\,869--872.


\bibitem{Shub}


Shubert H. {\em Kategorien} II. -- Berlin: Academie-Verlag, 1970. -- 148 p.


\bibitem{KR}


Кок А., Рейес Г.Э. {\em Доктрины в категорной логике}


// Справочная книга по матем. логике. Ч.\,I. Теория моделей.


-- М.: Наука, 1982. -- С.\,289--319.


\bibitem{BF}


Бордман Дж., Фогт Р. {\em Гомотопически инвариантные


алгебраические структуры на топологических пространствах}. --


М.: Мир, 1977. -- 408 с.


\bibitem{M}


Manes E.G. {\em Algebraic theories}. -- New~York: Springer,


1976. -- 356 p.


\bibitem{W1}


Wraith G.C. {\em Algebraic theories} // Lect. Notes Ser.


Math. inst. Aarhus univ. -- 1970. -- \No\,22. -- 131 p.


\bibitem{W2}


Wraith G.C. {\em Algebras over theories} //


Colloq. Math. -- 1971. -- V.\,23. -- \No\,2. -- P.\,181--190.


\bibitem{BCW}


Bass H., Connell E.H., Wright D.L. {\em Locally polynomial


algebras and symmetric algebras} // Invent. Math. -- 1977. --


V.\,38. -- \No\,2. -- P.\,279--299.


\bibitem{BA}


Бургин М.С., Артамонов В.А. {\em Некоторые свойства


подалгебр в многообразиях линейных $\Omega$-алгебр} //


Матем. сб. -- 1972. -- Т.\,87. -- \No\,1. -- С.\,67--82.


\bibitem{Bur}


Бургин М.С. {\em Свободные фактор-алгебры свободных


линейных $\Omega$-алгебр} // Матем. заметки. -- 1972. -- Т.\,11. --


\No\,5. -- С.\,537--544.


\bibitem{A4} 


Artamonov V.A. {\em The categories of free metabelian


groups and Lie algebras} // Comment. math. Univ. carolinae. --


1977. -- V.\,18. -- \No\,1. -- P.\,143--159.





\end{thebibliography} 


 


\vskip 1cm 


 


\post{\it Казанский государственный университет}{\it Поступила}


\post{}{31.08.1995} 


 


\end{document} 


